% Chapter 42: The Disasters That Confronted Classical Physics
\chapter{The Disasters That Confronted Classical Physics}

In a famous spring 1900 lecture, Lord Kelvin said physics' grand building was essentially complete, needing only finishing touches beneath two small clouds on the horizon. This was not arrogance: Newton explained apples through planets, Maxwell united electricity, magnetism, and light, and thermodynamics and statistics were unlocking heat. With two trophies polished bright, who would predict collapse?

Each cloud concealed a storm. One produced relativity, Chapters 38--40. The other, why hot things glow, toppled more people and produced quantum mechanics. We inventory its landfall here: in five years classical physics lost five battles in succession, each against inescapable experimental facts.

\step{A Counterintuitive Opening}

On a winter evening, look at a kitchen radiant electric cooktop. Just switched on, its surface already feels hot, though do not actually touch it, while remaining dark. Minutes later it glows dull red, then orange-red at higher power. An industrial furnace reaching higher temperatures turns bright yellow and even dazzling white.

Blacksmiths long made this a craft: one glance at heated iron, without a thermometer, reveals whether it is ready to forge. Notice the oddity hidden here. Put ceramic, charcoal, and firebrick beside that iron at equal temperature, and their glow has nearly the same color.

Color recognizes temperature, not material.

This is no artisan's illusion but a repeatable laboratory fact: anything at one thousand kelvin glows red, at fifteen hundred orange-yellow, and the solar surface at fifty-eight hundred bright white. Composition, hardness, and origin withdraw, leaving temperature alone directing color.

Now theory enters. Around 1900, physicists wielded their best tools, Maxwellian electromagnetism and statistical physics, to calculate how much a hot body emits in each color. The answer was neat, beautiful, and spectacularly wrong. Higher frequencies carried ever more energy, through ultraviolet, X-rays, and gamma rays, yielding infinite total energy. According to the best physics available, standing beside your oven should instantly fell you with lethal shortwave radiation.

Plainly it does not. Ovens run daily without sending people to hospital for that reason. Which assumption makes theory fail so completely? We will expose it, and in doing so explain two further oddities: fluorescent lamps shine white while merely warm, and you continually emit invisible light right now.

\step{Make a Guess First}

As usual, write intuitive answers on paper and grade yourself at the end.

First: iron and ceramic share a furnace at one thousand degrees. Is the iron redder, the ceramic redder, or are their colors essentially identical?

Second: illuminate a metal plate with intense red light for a long time or faint violet light briefly. Which more likely knocks out electrons? Or can both, with different delays?

Third: increasing furnace temperature takes the glow through red, orange, yellow, and white. Does the sequence end? Can enough heat turn it blue, and what follows blue?

Written them down? If sufficiently strong, prolonged red light seems enough, you agree with every physicist before 1905, and that is precisely the most decisively wrong answer.

\step{Hands-on Experiment}

\begin{experiment}[Separating Two White Lights with a Discarded CD]
Finding which colors light contains sounds expensive, but a discarded disc suffices.

\textbf{Procedure:} find a reflective, unlabeled CD or DVD. In a dark room, shine an incandescent or dimmable desk lamp obliquely onto its rainbow side. Bring your eyes near the disc and seek reflected light from different angles. At one angle a bright colored band appears. Record it, then repeat with a compact fluorescent lamp or a phone's plain white screen.

\textbf{What appears:} incandescent light gives a continuous rainbow, red through violet with smooth transitions. Fluorescent light breaks into isolated bright bands, especially striking green, separated by large gaps. The phone screen differs again: a narrow blue patch and broad yellow patch, with something missing between.

\textbf{Why:} over six hundred grooves per millimeter, about \(1.6\,\mu\mathrm{m}\) apart and comparable to light wavelengths, turn the disc into a vast diffraction grating. Grooves send different colors in different directions, as Chapter 35 explains, spreading white light's contents into an inventory. Incandescent light comes from heat and gives a continuous list; fluorescent gas discharge gives only a few entries.

\textbf{Think:} this explains nothing yet, merely juxtaposes facts. Why do hot sources provide continuous lists and cool ones sparse lists? Remember those bright bands: they will become coded letters from atoms.
\end{experiment}

\step{The Old Model Runs into Trouble}

Return to around 1900 and watch classical physics lose successive battles.

\textbf{Disaster one: thermal radiation totals infinity.}

First a fact: anything above absolute zero emits light, thermal radiation. At low temperatures it lies mainly in invisible infrared, leaving the newly heated cooktop dark. But a snake's tongue and a thermal imager can see it. You are radiating infrared throughout the room right now, a fact we revisit later.

The classical calculation seems entirely legitimate. Imagine cavity electromagnetic fields as countless oscillators: each frequency corresponds to a standing-wave mode, like innumerable differently tuned strings hanging in a room. Statistical physics in Chapter 31 offers tested equipartition: at thermal equilibrium each mode receives mean energy \(k_{\mathrm{B}}T\), with Boltzmann constant \(k_{\mathrm{B}}\) and temperature \(T\). Countless gas-molecule tests leave no reason for doubt.

Then disaster. Wavelengths can be arbitrarily short and frequencies arbitrarily high. Infinitely many high-frequency modes each demand \(k_{\mathrm{B}}T\); infinitely many nonzero amounts sum to infinity. The resulting Rayleigh--Jeans formula matches long wavelengths beautifully but goes wild at high frequency, predicting indefinitely rising intensity. Later generations named it the ultraviolet catastrophe.

The measured curve rises from low frequencies, peaks, then falls back toward zero. The Sun at about \(5800\,\mathrm{K}\) peaks near \(500\,\mathrm{nm}\), visible light. A thousand-kelvin filament peaks in infrared, with only its red tail reaching visibility, hence dull red. Your \(310\,\mathrm{K}\) body peaks near \(10\,\mu\mathrm{m}\) in far infrared. Theory is right at long wavelengths and utterly wrong at short ones, indicting a basic assumption rather than a coefficient.

\textbf{Disaster two: light knocks out electrons unreasonably.}

While verifying electromagnetic waves in 1887, Hertz noticed ultraviolet illumination helping sparks jump between metal electrodes. Over a decade later Lenard and others made precise experiments: light really ejects metal electrons, the photoelectric effect.

If light is a wave, then firmly established, three predictions seem indisputable:

\begin{itemize}[leftmargin=1.8em]
\item Stronger light carries more wave energy, so ejected electrons should move faster.
\item Any frequency, sufficiently intense and prolonged, should accumulate enough energy to eject electrons.
\item Weak light should produce a noticeable delay while electrons accumulate energy.
\end{itemize}

Experiments answer no three times. First, maximum electron energy depends only on color: violet beats blue, blue beats green. Increasing red intensity ten thousandfold produces neither more nor faster electrons. Second, each metal has a threshold frequency. Redder light ejects none even after a year under a searchlight; bluer light, almost invisibly weak, ejects them immediately. Third, sufficiently violet light produces no measurable delay, later bounded below a billionth of a second.

Pause over this direct rebuke to intuition: intense red loses to faint violet. Quantity loses to frequency. A classical delay estimate is worse still. Under ordinary illumination an atom-sized area receives roughly \(10^{-20}\,\mathrm{W}\), while ejecting an electron requires several electronvolts, order \(10^{-19}\,\mathrm{J}\). Accumulation should take minutes or hours. Experiment says under a billionth of a second, a billionfold mismatch.

\textbf{Disaster three: atoms should not exist.}

This is more complete. Fill a discharge tube with hydrogen, electrify it, and separate its glow with a prism or grating. Not a rainbow but specific bright lines appear. In 1885, school mathematics teacher Balmer found their wavelengths fitting a suspiciously simple empirical formula. Nobody knew why atoms glow or why only in those colors.

Stability was worse. Rutherford's scattering, detailed in Chapter 51, reveals a tiny heavy nucleus with outside electrons. The only available picture had orbiting electrons, accelerated charges that Maxwell's theory in Chapter 24 says must radiate and lose energy. Calculation spirals them into the nucleus in about \(10^{-10}\) seconds. Classical physics predicts collapse in one ten-billionth of a second: matter, including you, cannot exist.

All three point to one conclusion: classical calculations are not locally mistaken; the unspoken premise of energy flowing continuously in arbitrarily small portions fails microscopically.

\step{Inventing New Concepts}

Historically, repair began with something like cheating.

\textbf{Planck's patch, 1900.} Seeking the observed rise-and-fall curve, Planck tried an assumption uncomfortable even to himself: an oscillator at \(f\) exchanges energy not in arbitrary amounts but whole portions, each
\[
E = hf,
\]
where \(h = 6.63\times 10^{-34}\,\mathrm{J\cdot s}\) is a new natural constant, now Planck's constant. Energy resembles coins with a minimum denomination proportional to frequency.

Why does this cure the catastrophe? A high-frequency portion \(hf\) is costly while thermal motion offers only roughly \(k_{\mathrm{B}}T\) in pocket money. Beyond some frequency, even one portion is unaffordable and oscillators freeze out. Most of the infinitely many high-frequency modes receive nothing; the total becomes finite. The rising and falling curve matches experiment exactly.

Complete the analogy. A vending machine accepts whole coins; plentiful small change cannot buy an excessively priced item. This is not a partly closed tap: the flow is not slower, its minimum transaction larger. Planck disliked the patch for years as a mathematical device. A twenty-six-year-old patent clerk took it seriously.

\textbf{Einstein's photon, 1905.} Do not blame only oscillators, Einstein said; admit light itself behaves this way. At frequency \(f\), its exchanges with matter always come in \(hf\) packets, photons. Three photoelectric failures become three successes:

An electron swallows one photon at a time. The photon first buys the escape ticket, work function \(W\), priced differently by each metal. Leftover money becomes electron kinetic energy:
\[
hf = W + E_{\mathrm{k}}.
\]
Check each account. Insufficiently violet, \(hf < W\), cannot buy a ticket: greater intensity only queues billions of poor photons, and an electron serves one at a time. Thus the threshold defeats intensity. Adequate color transacts immediately upon arrival, with no saving period, hence no delay. Greater intensity brings more photons and more ejected electrons but leaves each electron's leftover energy unchanged. Intensity sets count, not speed. Every part fits.

\textbf{Bohr's energy levels, 1913.} Against discrete spectra and atomic collapse, Bohr likewise simply stipulated rules. Atoms allow only particular stationary energy states, which do not radiate, suspending Maxwell's accelerating-charge rule. Radiation occurs only upon a jump from a higher to a
lower level, with photon energy exactly their difference:
\[
hf = E_{\text{high}} - E_{\text{low}}.
\]
Discrete levels yield discrete differences and photon energies, hence bright lines. Hydrogen levels \(E_n = -13.6\,\mathrm{eV}/n^2\), with \(n = 1, 2, 3, \dots\), reproduce every mysterious Balmer wavelength. Immediately install a guardrail: Bohr electrons are not planets on definite nuclear routes; stationary states are not drawable trajectories. Orbiting electrons are leftover vocabulary, not the picture. Chapters 44 and 52 explain why the model partly works and where it ends.

\textbf{Photon momentum: Compton scattering, 1923.} If light comes in packets, they should carry momentum too. Chapter 40 gives zero-rest-mass particles \(p = E/c\), hence
\[
p = \frac{h}{\lambda}.
\]
Compton bombarded graphite electrons with X-rays. Classical waves would merely shake electrons, leaving scattered frequency unchanged. Instead, wavelength increased with scattering angle exactly as conservation predicts for a photon colliding with a stationary electron: about \(0.0024\,\mathrm{nm}\) at ninety degrees, twice that for backward scattering. The collision balances like two billiard balls. Photon momentum is measured, no longer metaphorical.

\textbf{Electron wavelength: de Broglie and diffraction, 1924--1927.} The last upheaval runs the other way. If waves can be counted as particles, de Broglie asked, might familiar particles be waves? He proposed that momentum \(p\) accompanies wavelength
\[
\lambda = \frac{h}{p}.
\]
Three years later Davisson and Germer sent electrons at nickel crystals while G. P. Thomson sent them through thin metal films. Both obtained alternating bright and dark diffraction patterns, waves' signature, impossible for particles alone. Electrons, electrically steered cathode-ray specks for thirty years, diffract.

The five battlefields are now inventoried. Do not conclude that light and electrons dress alternately as waves and particles, treating borrowed words as costumes. Neither classical category alone contains the facts. Photons and electrons each exist in one way; our old vocabulary is inadequate. Chapter 43 introduces a better one.

\step{The Model in One Sentence}

\begin{oneline}
Light and matter never exchange arbitrary fractions of energy but whole \(E = hf\) packets. Classical physics' disasters stem from the untested assumption that energy can be divided indefinitely.
\end{oneline}

\step{The Mathematical Translator}

With intuition ready, translate into calculable equations and test each against extremes immediately.

\textbf{First: \(E = hf\).} It gives packet energy for light at \(f\). Why this form? Proportionality is simplest, and the blackbody curve demands exactly it. Zero frequency gives \(E = 0\), no oscillation and no packet, sensible. Extreme frequency: gamma rays at \(f \sim 10^{20}\,\mathrm{Hz}\) carry about \(7\times10^{-14}\,\mathrm{J}\) per packet, roughly two hundred thousand visible photons together, hence lead shielding rather than glass. Low frequency: FM radio at \(f \sim 10^{8}\,\mathrm{Hz}\) carries about \(7\times10^{-26}\,\mathrm{J}\), too little for any receiver to count, explaining apparently continuous radio waves. Quantum behavior is averaged away because packets are tiny, not absent.

\textbf{Second: photoelectric \(hf = W + E_{\mathrm{k}}\).} It allocates a photon's energy after transfer to an electron. If \(hf\) equals \(W\), kinetic energy \(E_{\mathrm{k}} = 0\): the threshold frequency is \(f_0 = W/h\). If \(hf < W\), negative kinetic energy is impossible; no electron escapes and the equation announces its boundary. Doubling intensity changes no term and no individual electron energy, only transactions per second.

\textbf{Third: Balmer's formula.} Hydrogen's visible lines obey
\[
\frac{1}{\lambda} = R\left(\frac{1}{2^2} - \frac{1}{n^2}\right), \qquad n = 3, 4, 5, \dots
\]
with \(R \approx 1.097\times10^{7}\,\mathrm{m^{-1}}\). At \(n = 3\), \(\lambda \approx 656\,\mathrm{nm}\), hydrogen's famous red. As \(n \to \infty\), \(\lambda \to 365\,\mathrm{nm}\); increasingly crowded lines approach a series limit corresponding to escape from the atom. A purely empirical formula is fully explained by \(hf = E_{\text{high}} - E_{\text{low}}\), a receipt for energy levels, not coincidence.

\textbf{Fourth: de Broglie \(\lambda = h/p\).} Start with a sixty-kilogram person walking at one meter per second: \(\lambda = h/(mv) \approx 10^{-35}\,\mathrm{m}\), twenty orders smaller than a nucleus. You do not diffract through doorways because the wavelength is too small, not because the principle forbids it. An electron accelerated through one hundred volts instead has momentum about \(5.4\times10^{-24}\,\mathrm{kg\cdot m/s}\), giving \(\lambda \approx 0.12\,\mathrm{nm}\), comparable to crystal spacing. Crystals are natural electron gratings; the equation immediately explains their experimental use.

\textbf{A real-data estimate: photons from a lamp each second.} A \(100\,\mathrm{W}\) incandescent bulb converts only about five percent of electricity to visible light, the rest to heat, explaining its phaseout: about \(5\,\mathrm{J/s}\). With average visible wavelength \(550\,\mathrm{nm}\), one photon carries
\[
E = hf = \frac{hc}{\lambda} \approx \frac{6.63\times10^{-34}\times3.0\times10^{8}}{5.5\times10^{-7}} \approx 3.6\times10^{-19}\,\mathrm{J}.
\]
The count is about \(5 / 3.6\times10^{-19} \approx 1.4\times10^{19}\) per second. Ten quadrillion photons pour from your desk lamp each second. Continuous light means not unpackaged energy but packets too many and small for eyes or any instrument to distinguish their individual raindrops from the flood.

\begin{mathbridge}
\textbf{The catastrophe quantified, and Planck's resolution.}

Rayleigh--Jeans gives cavity radiation energy density near \(f\) as
\[
u(f) = \frac{8\pi f^2}{c^3}\,k_{\mathrm{B}}T.
\]
The \(f^2\) factor counts increasing standing-wave modes; \(k_{\mathrm{B}}T\) is equipartition. Integrating over all frequencies, \(\int_0^\infty u(f)\,df\), diverges. Catastrophe is no rhetoric but an infinity under an integral.

Planck replaces mean oscillator energy \(k_{\mathrm{B}}T\) with
\[
\bar{E} = \frac{hf}{e^{hf/k_{\mathrm{B}}T} - 1},
\]
yielding his distribution. Check both limits. At low frequencies, \(hf \ll k_{\mathrm{B}}T\), expansion \(e^{x} - 1 \approx x\) restores \(\bar{E} \approx k_{\mathrm{B}}T\), classical equipartition. At high frequencies, \(hf \gg k_{\mathrm{B}}T\), the denominator grows exponentially, suppressing \(\bar{E}\) toward zero, freezing modes and making the integral converge. Classical theory is not overturned but revealed as the \(hf \ll k_{\mathrm{B}}T\) approximation, explaining its two-century deception.

Finding the distribution's peak gives Wien's displacement law, \(\lambda_{\max} T \approx 2.9\times10^{-3}\,\mathrm{m\cdot K}\). Doubling temperature halves peak wavelength. Dull red, orange-red, yellow, and white all fit this equation.
\end{mathbridge}

\begin{challenge}
\textbf{Compton scattering: wavelength change from conservation.}

Treat scattering as an elastic collision. A photon carrying \(hc/\lambda\) and momentum \(h/\lambda\) strikes a stationary electron of mass \(m_{\mathrm{e}}\). The photon leaves at \(\theta\) with wavelength \(\lambda'\); the electron recoils with energy and momentum. Write energy conservation and momentum components in two directions, because momentum is a vector. Use relativistic electron energy \(E^2 = p^2c^2 + m_{\mathrm{e}}^2c^4\), eliminate electron energy and momentum, and obtain
\[
\lambda' - \lambda = \frac{h}{m_{\mathrm{e}}c}(1 - \cos\theta).
\]
At \(\theta = 0\), \(\cos\theta = 1\), so wavelength is unchanged: no collision, no change, sensible. At \(\theta = 180^\circ\), maximum shift \(2h/(m_{\mathrm{e}}c) \approx 4.9\times10^{-12}\,\mathrm{m}\) means maximum photon energy loss upon reversal. The constant \(h/(m_{\mathrm{e}}c) \approx 2.4\times10^{-12}\,\mathrm{m}\) is the electron Compton wavelength. Only photon energy \(hf\), momentum \(h/\lambda\), and conservation were used, both deriving the formula and directly evidencing photon momentum.
\end{challenge}

\step{The Model's Limits}

Powerful though these models are, each package states its scope.

First, photons are not miniature marbles. Particle behavior manifests in \(hf\) exchanges. No drawable ball trajectory belongs to a photon; propagation, interference, and obstacle-spreading still require waves, more deeply quantum field theory. A hard-ball picture immediately fails at double slits.

Second, duality is no transformation act. Photons and electrons never switch between wave and particle forms; the two classical categories are each incomplete. Experimental design reveals different phenomena, but the measured object remains one throughout. Chapter 43 nails this down with individual-particle double slits.

Third, in \(E = hf\), \(f\) is light's frequency as a periodic process, not a little ball shaking \(f\) times per second. The formula accounts for exchange, not internal construction.

Fourth, Bohr is transitional, not final. It precisely hits hydrogen lines but fails for helium and line intensities. Its achievement is introducing energy levels; Chapter 44's Schrödinger equation explains its inconsistency.

Fifth, most easily misused: quantum theory assigns classical territory, not wholesale rejection. For \(hf \ll k_{\mathrm{B}}T\), or momentum times size far exceeding \(h\), quantum formulas return to classical ones. Basketballs, planets, and microwave-oven waves still obey classical physics exactly. Common misuses: sufficiently strong light ejects any electron, false because frequency sets the threshold; color measures any source's temperature, false for fluorescent lamps and LEDs whose energy-level colors are temperature-independent, as shortly explained; quantum effects make the macroscopic world uncertain, false because your de Broglie wavelength is \(10^{-35}\) meters.

\step{Explaining the Opening Phenomenon}

Return to the kitchen stove and close each mystery.

Why does color recognize temperature rather than material? Thermal charge motion emits electromagnetic waves; equilibrium radiation follows Planck's distribution, whose only macroscopic parameter is temperature. Material emissivity changes quantity, not peak position. Equal-temperature iron, ceramic, and charcoal may differ in brightness but share essentially the same color. A blacksmith uses eyes as a rough spectrometer reading one Planck curve.

Why red, then yellow, then white? Wien says rising temperature proportionally shortens peak wavelength. At a thousand kelvin the peak lies deep in infrared, with only a red tail visible. At fifteen hundred, orange and yellow flood in. The Sun's fifty-eight hundred kelvin puts \(500\,\mathrm{nm}\) in the visible middle, all colors arriving together as white. Dimming an incandescent lamp reduces current, cools its filament, shifts the curve redward, and yellows or reddens the light. Your knob operates Wien's law by hand.

Why no ultraviolet catastrophe? High-frequency oscillators demand \(hf\) per portion, beyond thermal pocket money, and freeze out of the ledger. Ovens are quiet because Planck's constant is nonzero.

Why are fluorescent lamps white without being hot? They play a different game. Electrons in mercury-vapor discharge excite mercury atoms, which drop from higher to lower levels and emit specified lines according to \(hf = E_{\text{high}} - E_{\text{low}}\), strongest in ultraviolet. Wall phosphors absorb those photons and use their own level differences to emit broad visible spectra. The CD's discrete bands directly photograph mercury's energy differences. Those levels fix the color while the tube feels merely tens of degrees warm. Color thermometry fails because temperature never set that price list.

And your invisible light? Body temperature \(310\,\mathrm{K}\) gives peak wavelength around \(10\,\mu\mathrm{m}\) in far infrared. Each square meter of skin radiates several hundred watts while receiving environmental radiation, leaving much smaller net flow. This is industry, not metaphor: thermal imagers use detectors sensitive to \(10\,\mu\mathrm{m}\) photons, descendants of photoelectric semiconductor detectors, counting your infrared packets into pixels. Fire rescue, temperature screening, and night vision directly apply these formulas. Phone CMOS sensors, rooftop solar panels, and laboratory photomultipliers share the ledger: photon arrives, electron moves, one packet per entry.

\step{Thought Experiments and Challenges}

\begin{thought}[Turning Up Planck's Constant]
Quantum effects hide in daily life only because \(h\) is tiny. Imagine a cosmic knob gradually increasing \(h\). How far before you visibly see the world turn quantum?

Use a billiard ball, about \(0.2\,\mathrm{kg}\), rolling at \(1\,\mathrm{m/s}\) toward a five-centimeter pocket. Match wavelength to pocket width with \(\lambda = h/(mv)\), giving \(h = mv\lambda \approx 0.2\times1\times0.05 = 10^{-2}\,\mathrm{J\cdot s}\), roughly \(10^{32}\) times reality. There, balls spread and diffract toward pockets, entry becomes probabilistic, visible photons strike like bullets, lamps emit only a few each second, continuous light disappears, and hundred-meter race times become unmeasurable precisely.

The analogy resembles raising radio volume to hear a suppressed channel: quantum rules always play, merely too softly. It does not propose remodeling reality. Changing \(h\) also rewrites atoms, spectra, and chemical bonds, eliminating both billiard balls and you. The knob is imaginable, not operable.
\end{thought}

\begin{puzzles}
\item Two identical incandescent bulbs, one at full voltage shining white, the other undervolted and dull red: which emits more photons per second? \textbf{Hint:} compare more than power. Red photons cost less, buying more for equal energy. Imagine adjustable total powers, then compare how photon count and individual price change.
\item Green light just ejects a metal's electrons. Replace it first with double-intensity red, then half-intensity violet. What happens? \textbf{Hint:} is intensity in the photoelectric equation? Does it affect individual electron energy or transactions per second?
\item The solar peak lies in visible light and eyes are most sensitive there. Coincidence? \textbf{Hint:} reverse causation: in what illumination did eyes evolve? Which is cause and which effect? If the solar surface were ten thousand kelvin, what might earthly life see?
\item Why can electron microscopes see viruses while optical ones cannot? \textbf{Hint:} wavelengths limit resolution, blurring smaller details. Viruses are about \(100\,\mathrm{nm}\), visible light about \(500\,\mathrm{nm}\). Use \(\lambda = h/p\) for electrons accelerated through ten thousand volts and compare with a virus.
\end{puzzles}

The five-disaster inventory ends with a new key: \(hf\) energy packets and particles traveling with wavelength \(h/p\). We have only just put it in the lock. If an electron is a wave, what happens when one passes through both slits? That is stranger than disaster. See you in Chapter 43.
